\documentclass[10pt,a4paper]{amsart}
\usepackage[margin=19mm]{geometry}
\usepackage{amsmath,amssymb,mathtools}
\usepackage[scr=rsfs,cal=euler]{mathalfa}
\usepackage{xfrac}
\usepackage[unicode,hidelinks,bookmarks=false]{hyperref}
\newcommand{\ie}{i.e.\ }
\newcommand{\cf}{cf.\ }
\newcommand{\resp}{respectively\ }
\newcommand{\Subsection}{\subsection}
\providecommand{\nobreakdash}{\nobreak\hbox{-}}
\setlength{\emergencystretch}{2em}
\multlinegap0pt
\usepackage{bm}
\usepackage{bbm}
\usepackage{dsfont}
\usepackage{xspace}
\usepackage{cancel}
\usepackage{mathtools}
\usepackage{amsthm}
\makeatletter
\@ifundefined{thm}{\newtheorem{thm}{Thm}}{}
\@ifundefined{lemma}{\newtheorem{lemma}[thm]{Lemma}}{}
\makeatother
\newcommand{\al}{\alpha}
\newcommand{\la}{\lambda}
\newcommand{\p}{{\partial}}
\title[On a repulsion model with Coulomb interaction and nonlinear ]{On a repulsion model with Coulomb interaction and nonlinear mobility}
\author{Antonin Chodron de Courcel, Charles Elbar}
\date{Source: arXiv:2510.16894v1 (CC BY 4.0)}
\begin{document}
\maketitle

\noindent\emph{Source and licence.} On a repulsion model with Coulomb interaction and nonlinear mobility. Version arXiv:2510.16894v1 (CC BY 4.0), distributed under the Creative Commons Attribution 4.0 International licence, as recorded in the arXiv licence metadata for this exact version and shown on the abstract page https://arxiv.org/abs/2510.16894v1. Full rights notice: licenses/arxiv-2510.16894-CC-BY-4.0.txt.\par\smallskip

\noindent\textit{Editorial scope.} Counted unit: the Lemma whose source label is \texttt{lemma:supersolution} in the article (source0.tex lines 1576--1612, source label \texttt{lemma:supersolution}), delivered together with the article's own complete original proof of it (source0.tex lines 1633--1740, one \texttt{proof} environment of 108 lines). No mathematics is written, repaired or re-derived: statement and proof are the article's own text line for line. The proof is detached from the statement in the source: the article states the result at lines 1576--1612 and prints its proof at lines 1633--1740, and the proof environment's own header names the statement, so the association is the article's own. Required context retained uncounted: eq:primitive (lines 1307--1314). Every definition, display and label of those lines is retained unchanged. Omitted from this package and not redistributed: 1--1306, 1315--1575, 1613--1632, 1741--2125. Nothing inside a retained excerpt is omitted or altered: every retained body is delivered exactly as the article's lines are written, so no omission receipt of any kind accompanies this package. Citations: the retained excerpts carry no citation command at all, so no bibliography entry of the article is transcribed and no reference is redistributed. Reference adaptation: none was needed and none was made; every reference inside the delivered unit resolves inside it, so no unresolved reference or citation is produced. Printed slips kept literal: no printed word, symbol, hypothesis or number of the counted statement or of its proof was altered. Layout and class: the article's own class is \texttt{article} with packages including amsmath, amssymb, amsthm, hyperref, graphicx, color, comment, mathtools, cleveref, enumitem, xcolor, geometry, float, biblatex; the wrapper is the lane's pinned \texttt{amsart} editorial wrapper at 10pt on A4, which restates the theorem-like environments the retained text needs and adds the article's own macro definitions that the retained text needs (\texttt{\textbackslash al}, \texttt{\textbackslash la}, \texttt{\textbackslash p}), with extra wrapper packages (bm, bbm, dsfont, xspace, cancel, mathtools). The delivered numbering is the unit's own. Assets: no retained span contains a figure environment, a graphics-inclusion command, a PDF-page inclusion, a table or a TikZ picture; no image file is copied into this package, the package declares no asset dependency and no image file is redistributed with it. Licence: version arXiv:2510.16894v1 of this article is distributed under the Creative Commons Attribution 4.0 International licence, as recorded in the arXiv licence metadata for this exact version and shown on the abstract page https://arxiv.org/abs/2510.16894v1. A full rights notice is delivered at licenses/arxiv-2510.16894-CC-BY-4.0.txt. Characters: every retained excerpt is pure ASCII, so the delivered engine, XeLaTeX with its default Latin Modern text font, reports no missing character.
\begin{equation}\label{eq:primitive}
    \begin{cases}
    \p_t k + (\p_s k)^m_+ (k - s\bar u_0) = 0, & (t,s)\in (0,\infty)\times (0,1),\\
    k\vert_{s=0} = 0,\quad  
    k\vert_{s=1} = \bar u_0, \\
    k\vert_{t=0} = k_0.
\end{cases}    
\end{equation}

\begin{lemma}\label{lemma:supersolution}
    Let $C>0$ small enough, such that $C\bar u_0 \le 1$. Consider the map $ \sigma(u) := u^\frac{m}{m-1}$, and $\al \in (0,1)$ large enough, such that $2(1-\al)\sigma'(1)\le 1$. Let $(S_1,S_2,S_3)$ be respectively defined as
    \begin{align*}
        &S_1 := C\al \bar u_0,\\ 
        &\dot S_2 (t) = -(1-\al)^{m-1}\bar u_0^m \frac{S_3-S_2 + (1-\al)\sigma'(1) }{(S_3-S_2)^{m-1}}\sigma'(1)^{m-1},  \\
        &\dot S_3 (t) = (1-\al)^{m-1}\bar u_0^m \frac{1-S_3}{(S_3-S_2)^{m-1}}\sigma'(1)^{m-1},
    \end{align*}
    with initial condition $1\ge S_3^0> S_2^0> S_1$. Let us introduce the notation $\displaystyle u := \frac{s-S_2}{S_3-S_2}$. The function
    \begin{equation*}
     \tilde k_{C,S_2^0,S_3^0}(t,s) := \begin{cases}
        C^{-1}s & 0\le s \le S_1, \\
        \al \bar u_0 & S_1 < s \le S_2, \\
    \sigma(u) (1-\al )\bar u_0 + \al \bar u_0 & S_2 < s \le S_3, \\
    \bar u_0  & S_3 < s \le 1,
    \end{cases}
    \end{equation*}
    is a viscosity supersolution to \eqref{eq:primitive} on $[0,T_*]$, where we introduce
    \begin{align*}
        T_* &:= \inf\{t>0 : S_2(t) = S_1\},\\
         T^* &:=\inf\{t>0 : 2S_3(t) = 1 +S_3^0\},
    \end{align*}
    satisfying
    \begin{align*}
        T_* \ge C_m\bigg(\frac{S_2^0-S_1}{\bar u_0}\bigg)^m, \qquad
        T^* \ge C_m\bigg(\frac{1-S_3^0}{\bar u_0}\bigg)^m .
    \end{align*}
    for some constant $C_m>0$ depending only on $m$. 
    We finally record the following:
    \begin{equation*}
    \begin{cases}
      \forall t\in [0,T_*], &\quad S_2(t) \ge  S_2^0 - C_m\bar u_0 t^\frac{1}{m} , \\
      \forall t\in [0,\min(T^*,T_*)],&\quad  S_3(t) \ge S_3^0 + C_m' (1-\al)^{\frac{m-1}{m}} \bar u_0 t^\frac{1}{m}  \\
     \forall t\ge 0, &\quad S_3(t) \le S_3^0 + C_m'' \bar u_0 t^\frac{1}{m}.
    \end{cases}
    \end{equation*}
    
\end{lemma}

\begin{proof}[Proof of Lemma \ref{lemma:supersolution}]
Some details are skipped in the proof. A more detailed proof showing that a function (in this case a single and a double-vortex) is a viscosity solution can be found in the appendix.
By construction, $ \tilde k_{C,S_2^0,S_3^0}$ is continuous on $[0,T_*]\times [0,1]$, and $C^1$ outside of $\Gamma_1 \cup \Gamma_3$, where $\Gamma_i := \{(t,s) \in [0,T_*]\times [0,1] : s= S_i \}$. Consider $(t,s)$ such that $s< S_1$. We have
\begin{align*}
    &\p_t \tilde k_{C,S_2^0,S_3^0} (t,s) + (\p_s \tilde k_{C,S_2^0,S_3^0})^m_+(t,s) (\tilde k_{C,S_2^0,S_3^0}(t,s)-s\bar u_0) \\
    &= C^{-m}(C^{-1}-\bar u_0) s \ge 0.
\end{align*}
If $(t,s)$ is such that $s>S_3$ or $S_1<s<S_2$, the equation is trivially satisfied. Let us then consider $(t,s)$ such that $S_2 < s < S_3$. First,
\begin{align*}
    \p_t \tilde k_{C,S_2^0,S_3^0}(t,s) &= \frac{(1-\al)\bar u_0}{S_3-S_2} \sigma'(u) \big( -(1-u)\dot S_2 - u \dot S_3  \big) \\
    &= \bigg(\frac{(1-\al)\bar u_0}{S_3-S_2}\bigg)^m \bar u_0 \sigma'(1)^{m-1} \sigma'(u)  \big( (1-u) (S_3-S_2 + (1-\al)\sigma'(1)) - u(1-S_3)\big).
\end{align*}
Moreover,
\begin{align*}
    &\p_s \tilde k_{C,S_2^0,S_3^0}(t,s)^m(\tilde k_{C,S_2^0,S_3^0}(t,s) -s \bar u_0) \\
    &= \bigg(\frac{(1-\al)\bar u_0}{S_3-S_2}\bigg)^m \bar u_0 \sigma'(u)^m \big( (1-\al) \sigma(u) + \al - s \big) .
\end{align*}
Therefore, 
\begin{align*}
    &\p_t \tilde k_{C,S_2^0,S_3^0}(t,s) +\p_s \tilde k_{C,S_2^0,S_3^0}(t,s)^m(\tilde k_{C,S_2^0,S_3^0}(t,s) -s \bar u_0) \\
    &= \bigg(\frac{(1-\al)\bar u_0}{S_3-S_2}\bigg)^m \bar u_0\sigma'(u)\bigg( \sigma'(1)^{m-1}   \big( (1-u) (S_3-S_2 + (1-\al)\sigma'(1)) - u(1-S_3)\big) \\
    &+ \sigma'(u)^{m-1} \big( (1-\al) \sigma(u) + \al - s \big)  \bigg).
\end{align*}
The right-hand side above goes to zero as $u\to 1$ (hence $s\to S_3$). Moreover, the function 
\begin{multline*}
    u\mapsto \sigma'(1)^{m-1}   \big( (1-u) (S_3-S_2 + (1-\al)\sigma'(1)) - u(1-S_3)\big) \\
    + \sigma'(u)^{m-1} \big( (1-\al) \sigma(u) + \al - s \big)
\end{multline*}
is differentiable in $u$, and using  $\sigma'(u)^{m-1} = \sigma'(1)^{m-1} u$, $s= (S_3 - S_2) u + S_2$ its derivative is 
\begin{align*}
    &\sigma'(1)^{m-1}   \big( - (S_3-S_2 + (1-\al)\sigma'(1)) - (1-S_3)\big) 
    + \sigma'(1)^{m-1} \big( (1-\al) \sigma(u) + \al - s \big) \\
    &+ \sigma'(1)^{m-1}u \big( (1-\al) \sigma'(u) - S_3+S_2 \big),
\end{align*}
whose sign is that of
\begin{align*}
     &  - (S_3-S_2 + (1-\al)\sigma'(1)) - (1-S_3) 
    +  (1-\al) \sigma(u) + \al - s  \\
    &+ u \big( (1-\al) \sigma'(u) - S_3+S_2 \big)\\
    &\le  - (S_3-S_2 + (1-\al)\sigma'(1)) +S_3- s  
    + u \big( (1-\al) \sigma'(u) - S_3+S_2 \big) \\
    &\le  - (S_3-S_2 + (1-\al)\sigma'(1)) + S_3-S_2 + (1-\al)\sigma'(1) \\
    &= 0.
\end{align*}

Therefore, 
\begin{equation*}
    \p_t \tilde k_{C,S_2^0,S_3^0} + (\p_s \tilde k_{C,S_2^0,S_3^0})^m_+ (\tilde k_{C,S_2^0,S_3^0}-s\bar u_0) \ge 0,
\end{equation*}
for all $s\in (S_2, S_3)$.

We finally notice that
\begin{align*}
     D^{-}\tilde k_{C,S_2^0,S_3^0}(t,S_1) = \emptyset, 
    & \quad D^{-}\tilde k_{C,S_2^0,S_3^0}(t,S_3(t)) = \emptyset.
\end{align*}
Therefore, $\tilde k_{C,S_2^0,S_3^0}$ is a viscosity supersolution on $[0,T_*]$. Let us now derive a lower bound for $T_*$. 
First, we have $S_2(t)\le 1$ for all $t\ge 0$, so that
\begin{align*}
    \frac{d}{dt}(S_3-S_2) &= (1-\al)^{m-1}\bar u_0^m  \sigma'(1)^{m-1} \frac{1-S_2 + (1-\al)\sigma'(1)}{(S_3-S_2)^{m-1}} \\
    &\ge (1-\al)^{m}\bar u_0^m \frac{1 }{(S_3-S_2)^{m-1}}\sigma'(1)^{m}.
\end{align*}
Solving this equation, one obtains
\begin{equation}\label{ecartS3S2lowerbound}
    \forall t\ge 0, \quad S_3(t)-S_2(t) \ge C_m(1-\al) \bar u_0 t^\frac{1}{m}.
\end{equation}
We also record the following: if $2(1-\al)\sigma'(1)\le 1$ and $t\le T_*$ so that $S_2\ge S_1\ge0$,
\begin{align*}
    \frac{d}{dt}(S_3-S_2) 
    &\le C(1-\al)^{m-1}\bar u_0^m  \frac{1}{(S_3-S_2)^{m-1}}.
\end{align*}
Therefore,
\begin{equation}\label{ecartS3S2upperbound}
   \forall t\in [0,T_*],\quad  S_3(t)-S_2(t)\le  C_m'(1-\al)^{\frac{m-1}{m}}\bar u_0 t^\frac{1}{m}.
\end{equation}
For $t\in [0,T_*]$, we have $S_3-S_2 \le 1$. Using also \eqref{ecartS3S2lowerbound} and $2(1-\al)\sigma'(1)\le 1$,
\begin{align*}
    \dot S_2 &= -(1-\al)^{m-1}\bar u_0^m \frac{S_3-S_2 + (1-\al)\sigma'(1) }{(S_3-S_2)^{m-1}}\sigma'(1)^{m-1} \\
    &\ge -C_m \bar u_0 t^{-\frac{m-1}{m}}.
\end{align*}
Therefore,
\begin{equation*}
   \forall t\in [0,T_*],\quad  S_2(t) \ge S_2^0 -C_m \bar u_0 t^\frac{1}{m}.
\end{equation*}
This gives the lower bound on $T_*$. Finally, define for all $S_3^0 \le \lambda \le 1$
\begin{equation*}
    T_\lambda^* := \inf\{t>0 : S_3(t) = \lambda\}.
\end{equation*}
For all $t< \min(T_\lambda^*, T_*)$, we have using \eqref{ecartS3S2upperbound}
\begin{align*}
    \dot S_3(t) &\ge (1-\al)^{m-1}\bar u_0^m \frac{1-\lambda }{(S_3-S_2)^{m-1}}\sigma'(1)^{m-1}\\
    &\ge C_m (1-\la )(1-\al)^\frac{m-1}{m}\bar u_0 t^{-\frac{m-1}{m}}.
\end{align*}
Solving this equation gives, for all $t\le \min(T_\la^*, T_*)$,
\begin{equation*}
    S_3(t) \ge S_3^0 + C_m (1-\la )(1-\al)^\frac{m-1}{m}\bar u_0t^\frac{1}{m}.
\end{equation*}
Finally, using \eqref{ecartS3S2lowerbound},
\begin{align*}
    \dot S_3(t)& \le (1-\al)^{m-1}\bar u_0^m \frac{1 }{(S_3-S_2)^{m-1}}\sigma'(1)^{m-1} \\
    &\le C \bar u_0 t^{-\frac{m-1}{m}}.
\end{align*}
Thus, 
\begin{equation*}
   \forall t\ge 0,\quad  S_3(t) \le S_3^0 + C_m \bar u_0 t^\frac{1}{m}.
\end{equation*}
This gives a lower bound on $T_\la^*$. Taking $\displaystyle \la := \frac{1+S_3^0}{2}$ ends the proof.
\end{proof}

\end{document}
